\section{SELL-style rewriting and asynchronous adequacy}
\label{sec:sem}

The principal result of this section is Theorem~\ref{thm:sem-adequacy}:
an exact correspondence between asynchronous rewrite events and operational
peeling traces. A Solve event may occur while other checks still await values
assigned by earlier events. The stored rewrite residuals therefore cannot be
identified directly with the fully updated operational residuals. We relate
them through an invariant on pending incidences, then use propagation
normalization to obtain an exact endpoint that retains every log atom.

The two rules use persistent knowledge and linear residual resources. Their
semantics here is multiset rewriting; the logical notation presents an access
discipline rather than a SELL deductive calculus. The resulting trace can be
checked together with the residual success or failure witnesses of
Section~\ref{sec:certificates}. We first establish the local invariant and
normalization results, isolate the reusable proof principle, and then prove
the LDPC trace theorem in detail. Proposition~\ref{prop:sem-label-erasure}
subsequently determines which distinctions in the persistent-zone presentation
can be erased while preserving these executions.

\subsection{Resources, clauses, and execution labels}
\label{sec:sem-resources}

Fix a binary parity-check matrix $H$ and write
$\Nb(j)=\{i\in\Var:H_{j,i}=1\}$. A context is $M=(G,P,L)$: $G=H$ is fixed
static information, $P$ is a multiset of pairs $(i,b)$ representing persistent
facts $\mathsf{val}(i,b)$, and $L$ is a multiset of linear atoms
\[
 \mathsf{unk}(i),\quad \mathsf{pend}(j,i),\quad
 \mathsf{rem}(j,k,p),\quad \mathsf{step}(j,i,b).
\]
Here $i\in\Var$, $j\in\Chk$, $b,p\in\Ftwo$, and $k\in\mathbb N$.
A pending incidence is an edge whose variable has not yet been incorporated
into that check's stored residual. Its variable may already be known.
This distinction is necessary for asynchronous executions.

The two rules below act inside an arbitrary unchanged multiset context.
A premise in the persistent zone is read without being removed. Every displayed
linear premise is removed once, and every displayed conclusion is inserted once.
\begin{align*}
\textbf{Propagate:}\quad
 &!^K\mathsf{val}(i,b)\otimes\mathsf{pend}(j,i)
       \otimes\mathsf{rem}(j,k+1,p)
   \multimap\mathsf{rem}(j,k,p\xor b),\\
\textbf{Solve:}\quad
 &\mathsf{rem}(j,1,p)\otimes\mathsf{pend}(j,i)\otimes\mathsf{unk}(i)
   \multimap !^K\mathsf{val}(i,p)\otimes\mathsf{step}(j,i,p)
        \otimes\mathsf{rem}(j,0,0).
\end{align*}
The successor form of the propagation degree excludes underflow.
There are no rules that remove log atoms or create pending incidences.
The matrix is not read by either rule: the correspondence between pending
incidences and actual edges is established by the invariant below.

A rewrite event has label $\tau$ for Propagate and label $(j,i,p)$ for Solve.
An execution is a finite sequence of contexts and labelled events, with each
event justified by the applicable rule. Its trace is the ordered list of its
non-$\tau$ labels. We write
\[
 M\overset{T}{\Longrightarrow^{*}}M'
\]
when an execution from $M$ to $M'$ has exactly trace $T$.
The empty execution has trace $[\,]$; concatenation concatenates traces.
Thus the trace is part of the execution data, not an arbitrary ordering of the
final log multiset. In a proof assistant one may use computational derivation
objects or a proposition indexed explicitly by $T$; endpoint reachability alone
does not retain execution order.

For a trace $T$, let $\operatorname{Log}(T)$ be the multiset containing one
$\mathsf{step}(j,i,b)$ for each occurrence of $(j,i,b)$ in $T$.
Write $\uplus$ for multiset sum, so multiplicities are always retained.

\subsection{An invariant for delayed propagation}
\label{sec:sem-invariant}

\begin{defi}[Admissible asynchronous context]
\label{def:sem-admissible}
Let $K:\Var\rightharpoonup\Ftwo$ be a partial function and
$U=\Var\setminus\operatorname{dom}(K)$. For each $j$, let $Q_j$ be a finite
set and $p_j\in\Ftwo$. A context is admissible with data $(K,Q,p,T)$ if
$G=H$, its persistent zone is exactly
\[
 P=\{\!\{(i,K(i)):i\in\operatorname{dom}(K)\}\!\},
\]
and its linear zone is exactly
\begin{align*}
L={}&\{\!\{\mathsf{unk}(i):i\in U\}\!\}\uplus
 \{\!\{\mathsf{pend}(j,i):j\in\Chk,\ i\in Q_j\}\!\}\\
 &\uplus\{\!\{\mathsf{rem}(j,|Q_j|,p_j):j\in\Chk\}\!\}
 \uplus\operatorname{Log}(T),
\end{align*}
where
\begin{equation}
\label{eq:sem-pending}
 \Nb(j)\cap U\ \subseteq\ Q_j\ \subseteq\ \Nb(j),
 \qquad
 p_j=\bigxor_{i\in\Nb(j)\setminus Q_j}K(i).
\end{equation}
Every multiset comprehension above contributes each indicated atom exactly
once. The values in the XOR are defined because the first inclusion implies
$\Nb(j)\setminus Q_j\subseteq\operatorname{dom}(K)$.
\end{defi}

In particular, each variable has exactly one unknown token or exactly one
persistent value; conflicting values are excluded. Each row has exactly one
residual token, and pending incidences have multiplicity one. Membership-based
uniqueness would not suffice: two copies of the same residual atom have the
same parameters, but erasing one copy leaves the other behind.

\begin{defi}[Operational projection and normalized encoding]
\label{def:sem-projection}
The operational projection of an admissible context is
$\pi(M)=(K,U,s,d)$, where
\[
 d(j)=|\Nb(j)\cap U|,
 \qquad s(j)=\bigxor_{i\in\Nb(j)\cap\operatorname{dom}(K)}K(i).
\]
Let $D_j=Q_j\cap\operatorname{dom}(K)$ be its delayed known incidences.
The context is \emph{propagation-normalized} if every $D_j$ is empty.
For an operational state $\sigma$ satisfying the degree and syndrome invariants,
$\operatorname{Enc}_H(\sigma,T)$ is the admissible context with
$Q_j=\Nb(j)\cap U$, $p_j=s(j)$, and log multiset $\operatorname{Log}(T)$.
\end{defi}

The inclusion in~\eqref{eq:sem-pending} gives the disjoint union
$Q_j=(\Nb(j)\cap U)\mathbin{\dot\cup}D_j$. Consequently
\begin{equation}
\label{eq:sem-delayed}
 |Q_j|=d(j)+|D_j|,
 \qquad s(j)=p_j\xor\bigxor_{i\in D_j}K(i).
\end{equation}
These equations state precisely how the stored residual differs from the
operational residual. In a normalized context they agree, and
$M=\operatorname{Enc}_H(\pi(M),T)$. In an intermediate context the equalities
$|Q_j|=d(j)$ and $p_j=s(j)$ need not hold.

\begin{exa}[Two solves before propagation finishes]
\label{exa:sem-interleaving}
Let
\[
 H=\begin{pmatrix}1&0\\0&1\\1&1\end{pmatrix},
 \qquad y=(\star,\star).
\]
Initially $Q_1=\{1\}$, $Q_2=\{2\}$, $Q_3=\{1,2\}$, and all stored
syndromes are zero. Solve $(1,1,0)$ and then immediately solve $(2,2,0)$.
Both rewrites are enabled. The resulting assignment knows both variables,
but row~3 still contains $\mathsf{rem}(3,2,0)$ and both pending incidences.
Its operational degree is zero and its delayed set is $D_3=\{1,2\}$.
Either order of the two remaining propagations yields
$\mathsf{rem}(3,0,0)$ and leaves the trace
$[(1,1,0),(2,2,0)]$ unchanged. A proof that assumes all propagations follow
their originating solve immediately would not cover this execution.
\end{exa}

\begin{lem}[Initialization]
\label{lem:sem-init}
Let $\sigma_0$ be the initial peeling state for $y$. Then
$\operatorname{Enc}_H(\sigma_0,[\,])$ is admissible and normalized.
\end{lem}
\begin{proof}
Initialization sets $K$ to the visible coordinates of $y$,
$d(j)=|\Nb(j)\cap U|$, and $s(j)$ to the XOR of the known neighbors.
With $Q_j=\Nb(j)\cap U$, its complement inside $\Nb(j)$ is precisely the
known-neighbor set, so~\eqref{eq:sem-pending} holds with $p_j=s(j)$.
The encoding contains each required atom exactly once and no log atoms.
Also $Q_j\cap\operatorname{dom}(K)=\emptyset$.
\end{proof}

\begin{lem}[Propagation preserves the invariant and stutters]
\label{lem:sem-propagate}
Suppose $M$ is admissible with data $(K,Q,p,T)$ and a Propagate event at
$(j,i)$ with value $b$ yields $M'$. Then $M'$ is admissible with the same
$K$ and $T$, and $\pi(M')=\pi(M)$.
\end{lem}
\begin{proof}
The premises and uniqueness of the persistent value imply
$i\in Q_j\cap\operatorname{dom}(K)$ and $b=K(i)$. Replace $Q_j$ by
$Q_j\setminus\{i\}$ and $p_j$ by $p_j\xor b$; leave all other rows unchanged.
No unknown incidence is removed, so both inclusions
in~\eqref{eq:sem-pending} remain true. The set
$\Nb(j)\setminus Q_j$ gains exactly $i$, whose value is $b$; hence its XOR
changes by precisely $b$. Since $i\in Q_j$, its cardinality decreases by
one, as required by the new residual token. The unique pending and residual
tokens are replaced exactly as specified, and no other token or log changes.
The operational projection depends on $H$ and $K$, both unchanged.
\end{proof}

\begin{lem}[Solve preserves the invariant and simulates peeling]
\label{lem:sem-solve}
Suppose $M$ is admissible with data $(K,Q,p,T)$ and Solve $(j,i,b)$ yields
$M'$. Then $(j,i,b)$ is a valid operational peeling step from $\pi(M)$.
Moreover $M'$ is admissible with trace $T\mathbin{+\!+}[(j,i,b)]$, and
$\pi(M')$ is precisely the operational state after that step.
\end{lem}
\begin{proof}
The residual token gives $|Q_j|=1$ and $p_j=b$; the pending token gives
$i\in Q_j$; and the unknown token gives $i\in U$. Thus $Q_j=\{i\}$.
By~\eqref{eq:sem-pending}, $i$ is a genuine neighbor and every unknown
neighbor of $j$ lies in $\{i\}$. Therefore $\Nb(j)\cap U=\{i\}$.
Every other neighbor is known and outside $Q_j$, so
\[
 b=p_j=\bigxor_{v\in\Nb(j)\setminus\{i\}}K(v)=s(j).
\]
This is exactly the operational validity condition. No assumption is needed
about delayed incidences at other checks.

Set $K'=K\cup\{i\mapsto b\}$ and $U'=U\setminus\{i\}$.
At row $j$, put $Q'_j=\emptyset$ and $p'_j=0$. Its newly known-neighbor XOR
is $p_j\xor b=b\xor b=0$, agreeing with the rule. At any other row $\ell$
adjacent to $i$, the old inclusion $\Nb(\ell)\cap U\subseteq Q_\ell$
ensures $i\in Q_\ell$. It remains pending, so assigning its value does not
change the XOR over $\Nb(\ell)\setminus Q_\ell$. Retain $Q_\ell,p_\ell$.
At nonadjacent rows the sets and sums do not change. The pending inclusions
hold with $U'$ at every row.

Exactly one unknown token disappears and exactly one previously absent
persistent assignment is inserted. The forcing row's residual and pending
tokens are replaced as above; all other residual resources are retained.
One log atom is added. Hence the complete admissibility conditions hold.
Finally, fixing $i$ removes one unknown neighbor and adds $b$ to the
known-neighbor XOR at precisely the rows adjacent to $i$. This is the
operational update of $d$ and $s$, and proves the projection claim.
\end{proof}

\subsection{Normalization and exact trace adequacy}
\label{sec:sem-adequacy}

\begin{lem}[Propagation normalization and commutation]
\label{lem:sem-normalization}
Every admissible context $M$ with data $(K,Q,p,T)$ has a unique
propagation-normal form, denoted $\operatorname{NF}(M)$, and
\[
 \operatorname{NF}(M)=\operatorname{Enc}_H(\pi(M),T).
\]
Every maximal propagation-only execution has exactly
$\sum_j|D_j|$ events. Two distinct enabled propagations commute.
\end{lem}
\begin{proof}
A propagation is enabled exactly at a delayed incidence $(j,i)$ with
$i\in D_j$: its persistent value and pending token are present, and
$|Q_j|\geq1$. By Lemma~\ref{lem:sem-propagate}, such an event leaves $K$
fixed and removes exactly this incidence from $D_j$.
Thus the natural number $\sum_j|D_j|$ decreases by one at each event;
it is zero exactly when no propagation is enabled. Every maximal
propagation-only execution therefore terminates in the stated number of events.

At termination each pending set is $\Nb(j)\cap U$ by the two inclusions
in~\eqref{eq:sem-pending}. The invariant determines its syndrome as the XOR
of all known neighbors. Assignments, unknown tokens, static information,
and logs are unchanged. Exact multiplicities now force the terminal context
to be $\operatorname{Enc}_H(\pi(M),T)$, proving uniqueness.

For commutation, distinct enabled propagations at different rows modify
disjoint pending and residual tokens and read persistent values only.
At the same row they concern distinct pending variables, so the initial
pending degree is at least two and either remains enabled after the other.
Both orders remove those two incidences, decrease the degree by two, and
replace $p$ by $p\xor b\xor c$. Associativity and commutativity of XOR
make the results equal. The same incidence cannot be consumed twice because
its pending token has multiplicity one.
\end{proof}

\begin{prop}[Termination of rewriting]
\label{prop:sem-termination}
Every rewrite execution starting from an admissible context is finite.
Its number of events is at most the initial number $\sum_j|Q_j|$ of pending
incidences, and its number of Solve events is at most the initial $|U|$.
\end{prop}
\begin{proof}
Each rule consumes exactly one pending token and creates none. The preceding
lemmas preserve the invariant, so $\sum_j|Q_j|$ decreases by one at every
event and cannot become negative. Each Solve also consumes exactly one
unknown token; neither rule creates one.
\end{proof}

The confluence statement above concerns propagation with a fixed assignment.
It does not assert confluence of arbitrary Solve choices. On an inconsistent
input two different degree-one checks may force different values of the same
variable; both resulting traces can be locally valid, with inconsistency
remaining in another row. Certificate checking still tests the entire residual
system.

\begin{exa}[Why the endpoint includes a log]
\label{exa:sem-log}
For $H=[1]$ and $y=(\star)$, the valid step $(1,1,0)$ transforms
\[
 P=\emptyset,\quad
 L=\{\mathsf{unk}(1),\mathsf{pend}(1,1),\mathsf{rem}(1,1,0)\}
\]
into $P'=\{(1,0)\}$ and
$L'=\{\mathsf{rem}(1,0,0),\mathsf{step}(1,1,0)\}$.
An encoding of only the final peeling state would omit the step atom.
No further rule can erase it. Thus exact reachability of a log-free final
encoding is false even in this one-variable example.
The logged encoding $\operatorname{Enc}_H$ makes the correct endpoint explicit.
\end{exa}

\paragraph{An abstract adequacy principle.}
The preceding argument separates three observations of an execution: its
operational state, its ordered event trace, and the multiset of log atoms
retained at its endpoint. The following proposition isolates a sufficient
proof architecture. Its completeness hypothesis includes a local lifting
condition: forward simulation and unique silent normalization alone do not
ensure that an operational step can be performed by the concrete system.
The proposition is a mathematical abstraction of the argument, not an
additional mechanized result.

\begin{adequacyprinciple}
Let $\mathcal M$ be a concrete state space, $I\subseteq\mathcal M$ an
admissibility invariant, $\mathcal S$ an operational state space, and
$\mathcal A$ a set of observable events. Concrete transitions have labels
in $\mathcal A\cup\{\tau\}$, where $\tau\notin\mathcal A$ is silent;
operational transitions $\sigma\xrightarrow{a}_o\sigma'$ have labels
$a\in\mathcal A$. Let $\mathbb M(\mathcal A)$ be the finite multisets of
events, let $\pi:I\to\mathcal S$ and
$\ell:I\to\mathbb M(\mathcal A)$, and let
$\operatorname{Can}:\mathcal S\times\mathbb M(\mathcal A)\to I$.
Assume the following conditions.
\begin{enumerate}[label=(\arabic*),leftmargin=*]
\item \emph{Canonical representations.}
For every $\sigma,B$, the state $\operatorname{Can}(\sigma,B)$ is
silent-normal, with projection $\sigma$ and log $B$.
\item \emph{Invariant preservation and local simulation.}
Every concrete transition from a state in $I$ ends in $I$.
If $M\xrightarrow{\tau}_cN$, then $\pi(N)=\pi(M)$ and
$\ell(N)=\ell(M)$. If $M\xrightarrow{a}_cN$ for $a\in\mathcal A$, then
$\pi(M)\xrightarrow{a}_o\pi(N)$ and
$\ell(N)=\ell(M)\uplus\{\!\{a\}\!\}$.
These assertions concern transitions whose source lies in $I$.
\item \emph{Terminating canonical normalization.}
Silent reduction on $I$ terminates. Every silent-normal $N\in I$
satisfies $N=\operatorname{Can}(\pi(N),\ell(N))$.
\item \emph{Local lifting from canonical states.}
For every operational step $\sigma\xrightarrow{a}_o\sigma'$ and every
$B\in\mathbb M(\mathcal A)$, there exists a concrete transition
$\operatorname{Can}(\sigma,B)\xrightarrow{a}_cP$ with $\pi(P)=\sigma'$.
\end{enumerate}
Write $M\overset{T}{\Longrightarrow_c^*}N$ for a finite concrete
execution whose ordered nonsilent labels are $T$, and
$\sigma\overset{T}{\longrightarrow_o^*}\sigma'$ for an operational path
with exactly that event list. Then every $M\in I$ has the unique silent
normal form
\[
 \operatorname{NF}(M)=\operatorname{Can}(\pi(M),\ell(M)),
\]
and silent reduction from $I$ is confluent. Moreover:
\begin{enumerate}[label=(\roman*),leftmargin=*]
\item If $M\overset{T}{\Longrightarrow_c^*}N$, with $M\in I$, then
$N\in I$,
\[
 \pi(M)\overset{T}{\longrightarrow_o^*}\pi(N),\qquad
 \ell(N)=\ell(M)\uplus\operatorname{Log}(T),
\]
and $N$ silently normalizes to
$\operatorname{Can}(\pi(N),\ell(M)\uplus\operatorname{Log}(T))$.
\item For every $M\in I$ and every specified operational path
$\pi(M)\overset{T}{\longrightarrow_o^*}\sigma'$, there is a concrete
execution
\[
 M\overset{T}{\Longrightarrow_c^*}
 \operatorname{Can}(\sigma',\ell(M)\uplus\operatorname{Log}(T)).
\]
\end{enumerate}
Here $\operatorname{Log}(T)$ is the multiset image of the ordered list $T$;
the endpoint log need not determine that list. No determinism of operational
steps or confluence of the full concrete transition system is assumed.
\end{adequacyprinciple}
\begin{proof}
First consider silent reduction from $M\in I$. Termination gives a
silent-normal descendant: by well-founded induction on silent reduction,
take $M$ itself when it is normal, and otherwise take a silent successor
and apply the induction hypothesis to it. Invariant preservation keeps
every state in $I$. Repeated use of condition~(2) shows that a normal
descendant $N$ has projection $\pi(M)$ and log $\ell(M)$. Condition~(3)
therefore identifies it with $\operatorname{Can}(\pi(M),\ell(M))$.
This proves existence and uniqueness of the stated normal form. Any two
silent descendants of $M$ also have the same projection and log, and
each reduces to that same canonical state. Thus they have a common
silent descendant, proving confluence of silent reduction on $I$.

For (i), induct on the number of concrete events. For an empty execution,
the operational path is empty and the log is unchanged. For the induction
step, assume that the prefix ends at $P\in I$ with operational event list
$T_0$ and log $\ell(M)\uplus\operatorname{Log}(T_0)$. A silent next event
preserves both projection and log by condition~(2), so the same operational
path and event list suffice. An event labelled $a\in\mathcal A$ supplies
the operational step $\pi(P)\xrightarrow{a}_o\pi(N)$ and adds one occurrence
of $a$ to the log. Append this step to the path. Its ordered list is
$T_0\mathbin{+\!+}[a]$, whose multiset image is
$\operatorname{Log}(T_0)\uplus\{\!\{a\}\!\}$. In either case the
invariant is preserved. This proves the path and log assertions, and the
normalization result already proved gives the claimed endpoint.

For (ii), induct on the length of the specified operational path, allowing
any admissible initial concrete state. When the path is empty,
$\sigma'=\pi(M)$, and the silent normalization of $M$ supplies the required
execution. For a nonempty path, write its first step as
$\pi(M)\xrightarrow{a}_o\sigma_1$ and its remaining event list as $T_1$.
Normalize $M$ silently to $\operatorname{Can}(\pi(M),\ell(M))$.
Condition~(4) then gives an $a$-labelled transition to $P$ with
$\pi(P)=\sigma_1$. By condition~(2), $P\in I$ and
$\ell(P)=\ell(M)\uplus\{\!\{a\}\!\}$. Apply the induction hypothesis
to $P$ and the remaining operational path. Its endpoint is
\[
 \operatorname{Can}(\sigma',\ell(P)\uplus\operatorname{Log}(T_1))
 =\operatorname{Can}(\sigma',\ell(M)\uplus\operatorname{Log}(a::T_1)).
\]
Concatenating the initial silent normalization, the one observable event,
and the inductively obtained execution gives exactly the ordered list
$a::T_1$. This establishes (ii), including paths that stop before any
maximality condition holds.
\end{proof}

\paragraph{Instantiation to peeling.}
Fix $H$. Take $I$ to be the raw contexts of
Definition~\ref{def:sem-admissible}, $\mathcal S$ the coherent operational
states, $\pi$ the projection of Definition~\ref{def:sem-projection}, and
$\ell(M)$ the multiset of its $\mathsf{step}$ atoms, viewed as triples.
Silent transitions are Propagate; an observable transition is Solve with
its triple label. Define $\operatorname{Can}_H(\sigma,B)$ by the normalized
encoding with log multiset $B$. Equivalently, it is
$\operatorname{Enc}_H(\sigma,T)$ for any list with
$\operatorname{Log}(T)=B$: this is well-defined because that encoding
depends on $T$ only through its multiset image. Chronological order remains
part of the execution, not something recovered from $B$.

The normalized encoding supplies condition~(1).
Lemmas~\ref{lem:sem-propagate} and~\ref{lem:sem-solve} supply condition~(2),
and Lemma~\ref{lem:sem-normalization} supplies condition~(3), including
the exact canonical normal form. For condition~(4), let $(j,i,b)$ be any
valid operational step from $\sigma$. In its normalized encoding,
$Q_j=\Nb(j)\cap U=\{i\}$ and $p_j=s(j)=b$, so the unique residual,
pending, and unknown tokens enable precisely this Solve event, independently
of the previous log $B$. Lemma~\ref{lem:sem-solve} identifies the projection
of its successor with the prescribed operational successor. This verifies
local lifting without assuming the adequacy theorem being proved.

For the received-word initial state, Lemma~\ref{lem:sem-init} supplies
initialization. The other ingredients are silent propagation stuttering,
solving-step simulation, and terminating propagation normalization with a
unique normal form and commuting propagation steps. Together with the
explicit local lifting just verified, they give the following main theorem.
We retain its direct proof to display the exact assignments and logs in the
peeling instance.

\Needspace{20\baselineskip}
\begin{thm}[Main theorem: adequacy with exact trace labels]
\label{thm:sem-adequacy}
Let $\sigma_0$ be the initial state for $(H,y)$, and let
$M_0=\operatorname{Enc}_H(\sigma_0,[\,])$.
\begin{enumerate}
\item \emph{Soundness for arbitrary interleavings.}
If $M_0\overset{T}{\Longrightarrow^{*}}M$, then $T$ is a valid operational
peeling trace. Writing $\sigma_T$ for its replayed state, $M$ is admissible,
$\pi(M)=\sigma_T$, its log is exactly $\operatorname{Log}(T)$, and
\[
 \operatorname{NF}(M)=\operatorname{Enc}_H(\sigma_T,T).
\]
\item \emph{Completeness for each specified trace.}
For every valid operational peeling trace $T$ from $\sigma_0$,
\[
 M_0\overset{T}{\Longrightarrow^{*}}
       \operatorname{Enc}_H(\sigma_T,T).
\]
\end{enumerate}
These statements include empty and partial traces; neither requires that
peeling has already stalled.
\end{thm}
\begin{proof}
For soundness, induct on the events of the given execution, maintaining
admissibility, exact log equality, and equality between the operational
projection and replay of the trace prefix. The empty execution satisfies
these assertions by Lemma~\ref{lem:sem-init}. A Propagate event preserves
the invariant and projection by Lemma~\ref{lem:sem-propagate}; its label
$\tau$ does not extend the trace or the log. A Solve event is a valid
operational step by Lemma~\ref{lem:sem-solve}, and that lemma supplies
precisely the updated projection and appended trace/log entry. This proves
all assertions about the endpoint for every finite execution, including
those that end with unfinished propagations. Applying
Lemma~\ref{lem:sem-normalization} gives the stated normalized endpoint.

For completeness, induct on the length of the requested trace, maintaining
an execution to its exact logged normalized encoding. The empty trace is
realized by the empty execution. Suppose the prefix $T$ is realized at
$\operatorname{Enc}_H(\sigma_T,T)$ and its next valid step is $(j,i,b)$.
Since $\Nb(j)\cap U=\{i\}$ and $s(j)=b$, that encoding contains
$\mathsf{rem}(j,1,b)$, $\mathsf{pend}(j,i)$, and $\mathsf{unk}(i)$.
Apply Solve. For every other adjacent row $\ell$, the incidence $(\ell,i)$
was pending because $i$ was unknown; it is now a delayed known incidence.
No other delayed incidence is introduced, since the starting context was
normalized. Propagate these incidences in any order. Each remains enabled
until it is consumed; commutation and normalization follow from
Lemma~\ref{lem:sem-normalization}. By Lemma~\ref{lem:sem-solve} the projection
is the next operational state, and propagation does not change it.
The final context is therefore its exact normalized encoding with log
$\operatorname{Log}(T\mathbin{+\!+}[(j,i,b)])$.
The constructed block contributes exactly the requested triple, because
its remaining events are all $\tau$. Concatenating it with the prefix
execution establishes the induction step.
\end{proof}

The soundness direction applies to the events of the given execution, even
when a Solve occurs before earlier assignments have propagated to every
adjacent check. Completeness realizes each prescribed valid trace; it does
not replace that trace by another successful schedule. Normalization is
needed only to compare the stored residual resources with the fully updated
operational state. It changes neither the assignment nor the event history,
and the exact endpoint retains the log atoms that no rule removes.

The correspondence also links rewriting to certificate checking.
By Lemma~\ref{lem:core-residual}, replaying the emitted trace gives a residual
system whose solutions are exactly the compatible completions. Theorems
\ref{thm:core-success} and~\ref{thm:core-failure} then supply the success or
infeasibility contract for that same trace. This applies also to empty and
partial executions: certificate checking reconstructs the operational
residual from replay and does not require all concrete propagations to have
finished. The abstract proposition organizes the trace argument; the
certificate conclusions additionally use the binary linear algebra of
Section~\ref{sec:certificates}.

\begin{prop}[Maximal executions and stopping]
\label{prop:sem-stopping}
The endpoint of a maximal rewrite execution from $M_0$ is normalized.
Its projected unknown set is empty or is a nonempty stopping set.
\end{prop}
\begin{proof}
If a delayed known incidence remained, a propagation would be enabled,
contradicting maximality. In a normalized context $Q_j=\Nb(j)\cap U$.
If any check had exactly one unknown neighbor, its residual, pending, and
unknown tokens would enable Solve. Thus no residual degree is one.
When $U\neq\emptyset$, every check adjacent to $U$ has degree at least two,
which is precisely the stopping-set condition.
\end{proof}

This proposition concerns exhaustion of local peeling. In particular,
$U=\emptyset$ does not imply that all parity equations hold: a row may have
residual syndrome one. Both empty and nonempty residuals must undergo the
algebraic certificate checks.

\subsection{What the subexponential labels contribute}
\label{sec:sem-sell-role}

Following the resource distinction of linear logic and its subexponential
extension~\cite{Girard1987,NigamMiller2009SELL}, one possible logical
vocabulary uses the subexponential signature
$\Sigma=(\{S,K\},\preceq,\{S,K\})$, where $\preceq$ is the discrete
preorder and both labels admit weakening and contraction. The $S$ zone
contains static edge facts and reusable program clauses; the $K$ zone
contains learned value facts. Linear atoms have no such prefix.
The displayed clauses can be read as first-order multiplicative formulas,
with finite variable and check sorts, natural-number degree indices, and
$\Ftwo$-valued indices. Equivalently, use their ground instances, interpreting
successor and XOR in these index domains. This specifies the intended clause
syntax and structural permissions. In this paper the executable meaning
is exactly the two operational rules already defined. Structural proof rules
are not extra transitions of that operational system.

The subexponential presentation exposes the intended access discipline:
static data and reusable clauses occupy $S$, while newly learned values
occupy $K$. The following proposition identifies a distinction that is
semantically inessential for this transition system. It gives a controlled
comparison with one persistent zone and with ordinary multiset rewriting,
preserving the events used by the main theorem.

\begin{prop}[Erasing persistent-zone labels preserves these executions]
\label{prop:sem-label-erasure}
Consider the two-rule system above, with static edges represented by persistent
$\mathsf{edge}$ facts. Replace both $!^S$ and $!^K$ by one ordinary $!$,
retaining the atom predicates and event labels. On contexts in which edge
facts occupy $S$ and value facts occupy $K$, merging the two persistent zones
is an event-preserving strong bisimulation.
The same event-labelled transition system can be represented by ordinary
multiset rewriting, with the propagation rule
\[
\begin{split}
 \mathsf{val}(i,b)\otimes\mathsf{pend}(j,i)
             \otimes\mathsf{rem}(j,k+1,p)
 \multimap{}&\mathsf{val}(i,b)\otimes\mathsf{rem}(j,k,p\xor b),
\end{split}
\]
and Solve producing an ordinary $\mathsf{val}(i,p)$ token, while edge
tokens remain inert. This second correspondence holds on the images of
admissible contexts.
\end{prop}
\begin{proof}
For the first correspondence, merging is injective on the specified contexts:
the distinct predicate symbols $\mathsf{edge}$ and $\mathsf{val}$ recover
the original partition. Neither rule examines a persistent-zone label or a
preorder comparison. Propagate reads precisely the same value fact and
rewrites precisely the same two linear atoms before and after merging.
Solve consumes precisely the same three linear atoms and adds the same
value and log. Hence a transition in either system has a matching transition
with the same event label, and merging its endpoint gives the other endpoint.
There are only these two rules, so this proves both directions of the strong
bisimulation.

For ordinary rewriting, map each persistent edge and value fact to one
ordinary token and retain the linear multiset. Distinct atom predicates make
this map injective on admissible contexts. Propagate consumes the unique
value token and immediately recreates it unchanged; thus its net effect is
identical to persistent lookup. Solve introduces exactly the one new value
token, while the consumed unknown token and admissibility ensure that no
old value for that variable exists. The remaining consumed and produced
atoms coincide rule by rule, and neither system modifies an edge token.
Every applicable rule in either image therefore has the matching rule and
endpoint in the other image, with the same event label. Invariant
preservation ensures that subsequent contexts stay in these images.
\end{proof}

Proposition~\ref{prop:sem-label-erasure} separates the persistent-zone
vocabulary from the operational obligations used by certification. Its
translations retain the event labels, assignments, pending incidences,
residual tokens, and logs. Hence the invariant and exact trace correspondence
of Theorem~\ref{thm:sem-adequacy} survive erasure of the $S/K$ distinction.
The sufficient ingredients isolated by the abstract principle are invariant
preservation, projection and log behavior, canonical silent normalization,
and local lifting. Their use here does not depend on a comparison between
subexponential labels.

Differentiated persistent zones thus supply a specification discipline,
without an increase in operational expressiveness for these rules. The
proposition concerns the two stated rewrite systems, not an equivalence
between full logical calculi or a minimality theorem for resource-sensitive
representations. In particular, it does not establish adequacy for a focused
calculus of the kind studied in
\cite{Andreoli1992Focusing,NigamMiller2009SELL}. The certification result
remains the exact event-labelled correspondence despite delayed propagation,
with the log and normalization conditions made explicit.
